\section{Subresultants}

As before, for the entire section, let $f, g \in \mathbb{Z}[x]$ with $\deg f = n$ and $\deg g = m$. Moreover, we assume $m \leq n$. \issue{(Why is this assumption needed, and is it without loss of generality?)}

Resultant of polynomials is important, as it gives \issue{as} an efficient manner to understand if two polynomials have a common root or not. However, it does not give any other information but about gcd. Thus, we can extend the resultant definition into subresultants, which include all results of the Extended Euclidean Algorithm \cite{subresultant}. \issue{(What are the connections among subresultants, gcd, and the Extended Euclidean Algorithm?)}

First, we define Sylvester matrices for subresultants, similar to the resultant's matrix. Given $0 \leq k \leq m$,
\[
S_k(f,g)=
\begin{pmatrix}
f_n         &        &        &        & g_m         &        &        &        &        &        \\
f_{n-1}     & f_n    &        &        & g_{m-1}     & g_m    &        &        &        &        \\
\vdots      &        & \ddots &        & \vdots      &        & \ddots &        &        &        \\
f_{n-m+k+1} & \cdots & \cdots & f_n    & g_{k+1}     & \cdots & \cdots & g_m    &        &        \\
\vdots      &        &        & \vdots & \vdots      &        &        &        & \ddots &        \\
f_{k + 1}   & \cdots & \cdots & f_m    & g_{m-n+k+1} & \cdots & \cdots & \cdots & \cdots & g_m    \\
\vdots      &        &        & \vdots & \vdots      &        &        &        &        & \vdots \\
f_{2k-m+1}  & \cdots & \cdots & f_k    & g_{2k-n+1}  & \cdots & \cdots & \cdots & \cdots & g_k
\end{pmatrix}
\]
, which can be \issue{understand} as removing the last $k$\issue{-}columns of each submatrix of $f$ and $g$, and then removing the last $2k$ rows of the matrix.

Then, the  $k$-th subresultant of $f$ and $g$ is defined as \cite{subresultant}:
\[
\alpha_k = \det S_k(f, g)
\]
